% New results proved on 2026-10-15. Input after the main theorem statements.
% Requires the parent document's theorem and proof environments.
\section{A quantile lower bound under the same maintenance protocol}
\label{sec:quantile_extension}

The following result strengthens the expectation lower bound. Its statement
uses the same reporting event and the same cumulative unresolved area; it does
not introduce a separate deadline requirement. Write $\overline P$ for the
product prior on all candidate signs in the repeated-update family.

\begin{lemma}[Completion with a bounded candidate-sampling budget]
\label{lem:budget_completion}
Consider one batch of $m$ candidates in its validity window, under the
assumptions of the partial-recovery lemma. Let $\tau$ be the first round at
which at least $m/2$ candidates have determinate reports, or the end of the
window if this does not occur. Let $S$ denote that completion occurs, and let
$N(\tau)$ count only samples of this batch. Put
\[
 L=\frac{m\log(m/\delta)}{1024\kappa},\qquad
 q=\lfloor L/2\rfloor.
\]
If $L\ge4$, then
\[
 \overline P\{S,\ N(\tau)\le q\}<\frac12.
\]
\end{lemma}
\begin{proof}
Stop at the earliest of $\tau$, the window endpoint, and the first round with
$q+1$ candidate samples. At this bounded stopping time $\zeta$, the number of
candidate samples is at most $q+1<L$. On $\{S,N(\tau)\le q\}$, the budget
stopping condition has not fired before $\tau$, so $\zeta=\tau$ and at least
$m/2$ reports are determinate. If the displayed event had probability at least
$1/2$, the partial-recovery lemma at $\zeta$ would give
$\overline E N(\zeta)\ge L$, contradicting the deterministic bound
$N(\zeta)\le q+1<L$.

The extra sample in the stopping rule is needed. Stopping exactly at the
$q$th candidate sample would miss a completion that uses $q$ candidate
samples and later observations of the reference arm.
\end{proof}

\begin{theorem}[Regret quantiles required by rapid maintenance]
\label{thm:quantile_lower}
In the repeated-update family with $J$ validity windows, $M$ candidates per
window, $M\ge\max\{2,J\}$, $Jv\le\varepsilon/4$,
$d\le v/32$, $\eta=d/2$, and $0<\delta\le1/16$, suppose a policy is
$\delta$-reliable on the family and, for every sign configuration $\sigma$,
\[
 P_\sigma\{E_{\rm lab},\ \mathcal A_T(\eta)\le JMH/8\}\ge1-\delta.
\]
Then there is one fixed configuration $\sigma$ for which
\[
 P_\sigma\left\{R_T\ge
 2^{-20}\frac{\varepsilon JM}{d^2}\log\frac M\delta\right\}
 \ge\frac{11}{80}.
\]
This conclusion permits the policy to know the reference identities, the
arrival schedule, and the horizon.
\end{theorem}
\begin{proof}
Let $S_e$ denote completion of at least half of the current candidates before
the next reference arrives, and let $N_e$ be that batch's sample count at its
first such completion, or at the window endpoint when completion fails.
Set
\[
 L_0=\frac{M\log(M/\delta)}{32768d^2},\qquad
 q_0=\lfloor L_0/2\rfloor.
\]
The two directional Bernoulli divergences are at most $32d^2$, so the
partial-recovery lower bound is at least $L_0$. The assumptions imply
$d\le1/1024$ and $M\ge2$, hence $L_0>4$.
The preceding budget argument, with $q_0+1<L_0$, gives
\[
 \overline P(F_e)<\frac12,
 \qquad F_e:=\{S_e,N_e\le q_0\}.
\]
The argument is valid with all previous-window information included in the
history. Those observations and later reference distributions carry no
information about this batch's independent signs.

On the area-success event, any window without $S_e$ contributes at least
$MH/2$ area. Consequently $\sum_e1_{S_e}\ge3J/4$ on an event of prior
probability at least $1-\delta$. If $F=\sum_e1_{F_e}$, then
$\overline EF\le J/2$ and Markov's inequality gives
\[
 \overline P\{F\le5J/8\}\ge1/5.
\]
No independence across windows is asserted or needed. Intersecting these two
events has probability at least
$1/5-\delta\ge11/80$. On their intersection at least $J/8$ windows complete
with $N_e>q_0$, and therefore with $N_e>L_0/2$.
Each candidate observation in its birth window has gap at least
$\varepsilon-d\ge\varepsilon/2$. These are observations of distinct batches
in disjoint windows. Pathwise on that intersection,
\[
 R_T\ge\frac{\varepsilon}{2}\frac J8\frac{L_0}{2}
 =2^{-20}\frac{\varepsilon JM}{d^2}\log\frac M\delta.
\]
A mixture probability at least $11/80$ implies that one fixed configuration
has at least that probability. Thus no different worst environment is selected
for different windows.
\end{proof}

\begin{corollary}[Rapid maintenance with a joint regret quantile]
\label{cor:joint_quantile}
Let $\mathfrak R^{\rm joint}_\Theta(a)$ be the infimum of $r$ over policies
that are reliable on the full bounded model and satisfy
\[
 \inf_{\theta\in\Theta}
 P_\theta\{E_{\rm lab},\mathcal A_T(\eta)\le a,R_T\le r\}\ge1-\delta.
\]
For the repeated-update horizon
$H=128(M+J+1)\lceil32d^{-2}\log\{4K(K+1)/\delta\}\rceil$,
where $K=J(M+1)+1$, universal positive constants $c,C$ satisfy
\[
 c\frac{\varepsilon JM}{d^2}\log\frac M\delta
 \le\mathfrak R^{\rm joint}_\Theta(JMH/8)
 \le C\frac{\varepsilon JM}{d^2}\log\frac M\delta.
\]
\end{corollary}
\begin{proof}
For the lower bound, every policy feasible for the joint criterion satisfies
the area premise of Theorem~\ref{thm:quantile_lower}. Since
$11/80>\delta$, a joint regret bound strictly below its threshold is
impossible. Taking an infimum gives the claimed lower bound.
For the upper bound, the fine-checkpoint policy has the required area bound
on its simultaneous interval-coverage event, whose probability is at least
$1-\delta/2$. Independently of coverage, every available arm in the family has
gap at most $2\varepsilon$, so $R_T\le2\varepsilon T$ on every sample path.
The specified horizon and $J\le M$ yield the displayed upper order, with the
same universal constant as the rapid-maintenance upper bound.
\end{proof}

\section{A horizon-truncated lower bound across compatible task windows}
\label{sec:truncated_extension}

\begin{proposition}[Additive truncated area in one compatible experiment]
\label{prop:truncated_windows}
Consider finitely many batches indexed by $e$. Batch $e$ contains $m_e$
independent signs and has a deterministic validity window $W_e$ of length
$H_e$ that begins at that batch's arrival round. During that window each sign
is its arm's relative-goodness label and
its boundary distance exceeds $\eta$. The signs across all batches are
independent, other distributions do not encode them, and no information
about batch $e$ is available before its arrival. Samples remain interleaved
and adaptive. Suppose every sign configuration has familywise reporting error
at most $0<\delta\le1/16$. Let $0<\kappa_e<\infty$ bound the two directional
divergences in batch $e$, and define
\[
 h_e=\min\left\{H_e,
 \left\lceil\frac{m_e\log(m_e/\delta)}{1024\kappa_e}\right\rceil-1\right\},
 \qquad S=\sum_e m_eh_e.
\]
When $S>0$, one fixed full sign configuration satisfies
\[
 P\{\mathcal A_T(\eta)\ge S/8\}\ge1/7,
\]
provided every selected window lies before $T$ and every candidate--time pair
is counted only once. The windows may overlap when they concern distinct
arms and their stated qualification structure is jointly realizable.
\end{proposition}
\begin{proof}
At the $r$th round after batch $e$ arrives, $r\le h_e$, the entire experiment
has supplied at most $r$ samples of that batch. The contrapositive of the
partial-recovery lemma shows that at least half the batch is determinate with
probability less than $1/2$. Its expected unresolved count is therefore at
least $m_e/4$. Let $Y$ sum unresolved candidate--time pairs over the first
$h_e$ rounds of all selected windows. All these pairs have boundary distance
larger than $\eta$, so $Y\le\mathcal A_T(\eta)$, while $0\le Y\le S$ and
$\overline EY\ge S/4$. If $p=\overline P(Y\ge S/8)$, then
$S/4\le S/8+(7S/8)p$, giving $p\ge1/7$. Select a fixed configuration from
the common product prior. Compatibility of the windows is essential: this
proof cannot sum bounds from mutually incompatible reference environments.
\end{proof}
